\chapter{Background}\label{ch:background}

The following chapter is intended to give readers an overview about the
subject and state of the art of open source backup solutions.
It will also define important concepts used throughout this thesis.

\section{Current Developments and Challenges}\label{sec:current-developments}

Since the need for recovering data after a catastrophic failure such as a hardware
malfunction is not new there are a number of existing solutions to this problem.

\paragraph{rsync} A possible way of implementing such a solution is based on the
rsync tool which allows copying directories while replacing files existing in
other directories with hard links to them~\cite[argument `--link-dest']{rsync}.
This in turn allows saving storage space when performing periodic backups and a
backup contains unchanged files which are already contained in a previous backup.
The operating system ensures that when deleting a backup the files being referenced
by hard links from other backups are not actually deleted but remain accessible.

Examples of this approach are Back In Time~\cite{backintime} and rsnapshot~\cite{rsnapshot}.

\paragraph{Full/incremental backups} Another approach consists of periodically performing
backups without any data deduplication, called `full' backups.
Backups happening between them are stored as changes which happened since the last backup,
being called `incremental' backups as a result~\cite[section `OPERATION AND DATA FORMATS']{Duplicity}.

This allows to apply transformations such as encryption and compression to the generated
backups~\cite[section `DESCRIPTION']{Duplicity}, which are not possible when performing
backups based purely on rsync since it was originally designed to be
`a fast, versatile, remote (and local) file-copying tool'~\cite{rsync}.
Another advantage is the ability to detect and store differences between files of different
backups, allowing for storage space savings where rsync would require creating a new
file containing only slight differences from a previous backup~\cite[section `OPERATION AND DATA FORMATS']{Duplicity}.

A drawback of this approach is the need to apply all incremental backups performed since
the last full backup to it when restoring a single backup~\cite[section `OPERATION AND DATA FORMATS']{Duplicity},
requiring users to balance the time required for backup restoration with the processing
and storage costs associated with creating and keeping periodic full backups.
It also requires keeping backups which are no longer of interest to allow for
incremental backups based on them to be restored.

Examples of this approach are Déjà Dup~\cite{DejaDup} which is based on Duplicity~\cite{Duplicity}.

\paragraph{Content-addressable storage} An alternative to full/incremental backups which
still allows applying transformations to the generated backups while avoiding the mentioned
problems of backups depending on each other consists of splitting the data to be backed up into chunks.
These chunks are stored in a key-value store based on a key derived from their
content, allowing different backups which share common data to save storage space
by referencing the chunks containing this data multiple times~\cite[section `Internals']{Borg1}.
When restoring backups the chunks containing its data are retrieved from the key-value store
and the splitting operation is reversed, thereby avoiding the need to access other backups.

The use of a shared key-value store also enables backups to be performed in parallel
while allowing deduplication of data between them, which is not possible when using
Duplicity~\cite[question 284867]{DuplicityLaunchpad} but also dependent on the
implementation used.
An example of this is Borg 1 requiring exclusive access~\cite[section `Frequently asked questions']{Borg1}
while Borg 2~\cite[section `Important notes 2.x']{Borg2} and restic~\cite[Issue 4321]{resticGit}
allow multiple backups to the same key-value store to be performed in parallel.

A drawback of this approach is that when a backup is deleted its chunks remain in the
key-value store, continuing to occupy storage space even if not referenced by any remaining backups.
To reclaim the storage space used by such chunks an operation has to be run periodically which
calculates the chunks still referenced and deletes the rest~\cite[section `Usage']{Borg1}.
The cost of performing this operation is increased by many implementations such as
Borg 1~\cite[section `Frequently asked questions']{Borg1},
Borg 2~\cite[section `Important notes 2.x']{Borg2} and
Restic~\cite[section `Removing backup snapshots']{restic} requiring exclusive access
to the key-value store for its duration, preventing backup creation or restoration.

Another problem is the need to balance the amount of created chunks with the
efficiency of deduplication.
Increasing the number of chunks by reducing their individual sizes also increases
the resources required to manage them while simultaneously allowing for a more
fine-grained deduplication~\cite[section `Additional Notes']{Borg1}.
To reduce the impact of the number of chunks on performance some programs
like Borg 1~\cite[section `Internals']{Borg1} and restic~\cite[section `Terminology']{restic}
combine multiple chunks when storing them, for example as a single file in the file system.

Examples of this approach are Borg 1~\cite{Borg1}, Borg 2~\cite{Borg2}, restic~\cite{restic},
Duplicacy~\cite{Duplicacy}, Duplicati~\cite{Duplicati}, Kopia~\cite{Kopia} and
Pika Backup~\cite{Pika} which is based on Borg 1.

\paragraph{Other} Besides the aforementioned solutions there are a number of alternatives
like Bacula~\cite{Bacula} which for the purpose of clarity will not be further investigated.

\section{Lock-Free Deduplication}

\subsection{Duplicacy}

In their 2020 paper~\cite{Duplicacy} Gilbert Chen et al.\ proposed a solution to
the problem of periodic chunk deletion operations requiring exclusive access when
using a shared content-addressable storage for backup operations.
Their solution consists of a new backup program called `Duplicacy' which, while
using the aforementioned content-addressable storage approach, allows backup
and chunk deletion operations to run in parallel without relying on exclusive
access to the key-value store to ensure data consistency.

Duplicacy allows performing the creation of chunks either by using a fixed chunk size
or are more flexible approach called
`variable-size chunking'~\cite[section `Incremental Backup and Full Snapshot']{Duplicacy}.
While using a fixed chunk size has a smaller computational overhead in comparison
to the variable-size chunking algorithm it will fail to deduplicate data which was
moved, for example by inserting data into an existing file.
Variable-size chunking overcomes this problem by using a rolling hash function at the
cost of an increased need for computational resources.
This rolling hash function can be efficiently computed as a sliding window over the data
and is used to identify boundaries between chunks, allowing them to be recognized
after a move~\cite[section `Chunking']{Duplicacy}.

While Duplicacy does not combine chunks it uses a different technique to reduce the
number of objects stored.
When creating a backup consisting of multiple files they are first sorted and
combined into a single data stream, an approach similar to the well-known
`tar' archival program.
The resulting stream is then split into chunks, making the number of chunks
only dependent on the amount of data~\cite[section `Naming Chunks']{Duplicacy}
while some other programs such as Borg 1 create at least as many chunks as files
present in the backup~\cite[section `Internals']{Borg1}.

Besides these features Duplicacy allows applying compression and encryption to
the generated backups~\cite[section `Compression and Encryption']{Duplicacy},
which will not be in the focus of this thesis.

\subsection{Algorithm}\label{sec:algorithm}

To avoid the need for exclusive access Duplicacy uses an algorithm called
`Lock-Free Deduplication' which operates on a central key-value store,
called the `(file) storage' in their paper but referred to as the `\textit{repository}'
throughout this thesis to better align with the terminology used in
programs such as Borg~\cite[section `Internals']{Borg1, Borg2} and
restic~\cite[section `Terminology']{restic}.
The repository stores chunks as separate objects with a name based on their
individual contents, avoiding the need for a central database to track which
objects contain which chunks and the associated problems of parallel
access~\cite[section `Naming Chunks Based Backup']{Duplicacy}.

It is used by multiple actors which can perform operations such as creating
and deleting backups on the repository, called `\textit{clients}'.
To create a backup a client generates chunks as outlined in the previous
section but when uploading them to the repository follows the following policy:

\begin{samepage}
    \paragraph{Policy 1}\label{para:policy1} (Checking before Uploading).
    \textit{
        Before uploading a new chunk, the backup procedure always performs
        a lookup with a filename determined only by the chunk content hash.
    }~\cite{Duplicacy} \\
\end{samepage}


Since this will identify existing chunks even when created by other clients
their backups will share identical
chunks+\cite[section `Naming Chunks Based Backup']{Duplicacy}.
The backup itself consists of a `\textit{manifest file}' which contains the
necessary metadata to restore it with the chunks uploaded to the repository.
While this manifest file does not store the data of the backup directly it
can still reach a large size should many chunks be referenced.
To reduce the storage space it occupies it is also split into chunks and
stored in the repository, reducing the number of chunks directly referenced
by it~\cite[section `Naming Chunks']{Duplicacy}.
\Cref{fig:repository-structure} shows the resulting schematic structure with
chunks produced by backup data in green and the ones produced by manifest files
in red.

\begin{figure}
    \includesvg[width=0.75\textwidth, pretex=\small]{repository_structure.drawio.svg}
    \caption{Schematic structure of a repository.}\label{fig:repository-structure}
\end{figure}

While the paper does allow replacing chunks to allow for a lock-free backup
operation~\cite[Theorem 1]{Duplicacy} it does not state whether this is the case
for manifest files, but upon request this turned out to be not the
case~\cite[topic `Duplicacy paper data loss ']{Duplicacy-Forum}.

When deleting a backup by simply deleting chunks which would become unreferenced
there is the possibility of a chunk being deleted while another client creates a
manifest file referencing it, causing the restoration of the associated backup to
fail.
To prevent this instead of deleting them such chunks are first turned into `\textit{fossils}'
by renaming them, allowing them to be referenced by the following policy:

\begin{samepage}
    \paragraph{Policy 2}\label{para:policy2} (Fossil Accessing).
    \textit{The following two rules present the access policy for fossils.}
    \begin{enumerate}
        \item \textit{
            A restore, list or check procedure that reads existing backups can
            read the fossil if the original chunk cannot be found.
        }
        \item \textit{
            A backup procedure cannot access any fossils. That is, it must upload
            a chunk if it cannot find the chunk, even if a corresponding fossil exists.
        }
    \end{enumerate}~\cite{Duplicacy} \\
\end{samepage}

This step is called the `\textit{fossil collection step}' and consists of downloading
all existing manifest files from the repository, using them and a list of manifest files
of backups to be deleted to calculate the set of unreferenced chunks and turn them into
fossils~\cite[section `Two Step Fossil Deletion']{Duplicacy}.

The actual deletion of the created fossils is performed by the `\textit{fossil deletion step}',
which will be delayed until the fossils will no longer be referenced by any manifest files
which are in the process of being created as by the following policy:

\begin{samepage}
    \paragraph{Policy 3}\label{para:policy3} (Fossil Deleting).
    \textit{The following two conditions should be met before deleting fossils permanently.}
    \begin{enumerate}
        \item \textit{
            For each backup client, there is a new backup that was not seen by
            the fossil collection step.
        }
        \item \textit{
            For each backup client, the new backup must have finished after
            the fossil collection step.
        }
    \end{enumerate}~\cite{Duplicacy} \\
\end{samepage}

Should these two conditions be met downloading all manifest files existing in the
repository will find any new manifest files which may still reference the fossils.
Should a fossil be referenced it is turned back into a chunk (called `\textit{fossil recovery}'),
otherwise it is permanently deleted.
The data required for these operations is taken from a `\textit{fossil collection file}'
which was created by the fossil collection step and stored outside the repository
locally by the client performing these operations~\cite[section `Two Step Fossil Deletion']{Duplicacy}. 

\subsection{Race condition}\label{sec:race-condition}

While the paper claims that the aforementioned algorithm is lock-free~\cite[Theorem 2]{Duplicacy}
there is actually a possibility of data-loss when multiple clients perform the fossil
collection and deletion steps at the same time, as shown in \Cref{fig:race-candition}.

\begin{figure}[H]
    \includesvg[width=0.95\textwidth, pretex=\small]{fossil_race_condition.drawio.svg}
    \caption{Example of parallel execution causing data loss.}\label{fig:race-candition}
\end{figure}

In the shown example there exist 3 clients which operate on a central repository.
It begins with a client creating a backup called `\#1', which consists of uploading
a new chunk and a manifest file referencing it.
After that another client performs the fossil collection step required to delete
this backup, which deletes its manifest file and turns the new chunk into a fossil.
Next, every client creates a backup with its associated manifest file to allow the
fossil deletion step to proceed under \nameref{para:policy3}.

However, before it is executed, another client performs the fossil collection step
required to delete on of these backups, called `\#2'.
Should it have included re-uploading the same chunk which was turned into a fossil
earlier the fossil collection step will turn it into a fossil should none of
the other manifests reference it.
Since \nameref{para:policy2} implies that there can only exist at most one fossil per chunk
name the renaming operation will replace the existing fossil.

This creates a dangerous situation where multiple fossil collections contain the same
fossil but have different timestamps associated with them when checking \nameref{para:policy3}.
Should the fossil deletion step be executed it can't distinguish between fossils created
by its fossil collection step and fossils of the same chunks created by different fossil
collection steps.
This in turn can cause it to delete them even before \nameref{para:policy3} should allow the
fossil deletion step associated with the fossil collection step that created to be executed,
violating the assumption that the fossil will no longer be referenced by manifest files
which are in the process of being created.

In the example this is exploited by a client creating a manifest file called `\#3` referencing
the chunk which was turned into a fossil but looked up before and therefore not re-uploaded
again.
Should this manifest file be created after the fossil deletion step downloaded all manifest
files it will delete the fossil and therefore cause restoration of the associated backup to fail.

\subsection{Solution}

To prevent the race condition explained in \Cref{sec:race-condition} the parallel execution
of the fossil collection and deletion steps has to be constrained.

\begin{samepage}
    \begin{thm}\label{thm:singleclient}
        Assuming the manifest files of backups to be deleted are removed during the
        fossil collection step.
        By limiting fossil collection and deletion steps to a single client and the
        number of pending fossil deletion steps to 1 fossil collection and deletion
        steps influencing each other does not disrupt the Lock-Free Deduplication algorithm.
    \end{thm}
\end{samepage}

\paragraph{Proof} The following interactions are possible:
\begin{enumerate}
    \item A fossil collection step reverts a fossil recovery by turning the chunk back into a fossil

        Should it revert a fossil recovery the fossil was once determined to be
        referenced by a fossil deletion step and should therefore only be fossilized
        after the referencing manifest files seen by the fossil deletion step were
        either removed or are about to be deleted by the fossil collection step.
        Since the fossil deletion step completed before the fossil collection step
        started any remaining manifest files seen by it are also seen by the fossil
        collection step.

    \item A fossil collection step creates a fossil operated on by a fossil deletion step

        Should it create a fossil on which a fossil deletion step operates the
        fossil deletion step may recover or delete the fossil which will be proven
        not to disrupt the Lock-Free Deduplication algorithm by case 3 and 4.

    \item A fossil deletion step recovers a fossil

        Should the fossil deletion step recover a fossil it was once thought to be
        unreferenced and should therefore only be recovered if it is referenced again.
        Since a fossil deletion step prevents fossil collections steps from deleting
        manifest files during its execution a seen manifest file will still exist at
        the time the fossil is recovered.

    \item A fossil deletion step deletes a fossil

        Should the fossil deletion step delete a fossil it should only do so when
        there is no manifest file still referencing it.
        The fossil deletion step can only receive fossils from its fossil
        collection step or one which was executed before it and can only start after
        \nameref{para:policy3} ensures no new manifest files can reference fossils
        created before its fossil collection step ended.
        Since there can be at most one pending fossil deletion step the fossil
        deletion step can only operate on fossils created by its associated fossil
        collection step and therefore any manifest file still referencing the fossil
        will be seen by the fossil deletion step.

    \item A fossil deletion step creates a chunk operated on by a fossil collection step

        Should a fossil deletion step recover a fossil into a chunk on which a
        fossil collection step operates the fossil collection may revert it which
        was proven not to disrupt the Lock-Free Deduplication algorithm by case 1.
\end{enumerate}

While \Cref{thm:singleclient} might seem like a regression to requiring exclusive access to
the repository during the fossil collection and deletion steps it only applies to these operations.
Since the Lock-Free Deduplication algorithm still guards against clients creating backups
in parallel they can continue to do so while one of them performs the fossil collection
or deletion step.
It seems that these constraints are acceptable for most use cases since Duplicacy
itself~\cite[page `prune']{Duplicacy-Wiki} uses a similar approach, but whether this race
condition is known to them or it is done because
`This can greatly simplify the implementation'~\cite[topic `Prune command details']{Duplicacy-Forum}
could not be determined.
